(1) 弦方程为$\rho_\ell y_{tt}=T(t)y_{xx}$，固定端本征函数$\sin(n\pi x/L)$，$n=1,2,\ldots$。写$y=\sum_nq_n(t)\sin(n\pi x/L)$，正交投影得$\ddot q_n+2\beta_n\dot q_n+\omega_n^2[1+h\cos(\Omega t)]q_n=0$，其中$\omega_n=(n\pi/L)\sqrt{T_0/\rho_\ell}$。空间模态不混合是张力沿全弦均匀的前提。
(2) 在$\Omega=2\omega_n$附近，取$q=A(t)\cos(\Omega t/2)+B(t)\sin(\Omega t/2)$，幅变量相对快振动缓慢。在$|h|\ll1$且先忽略阻尼时，对一周期平均得到二维慢变矩阵，其共振耦合幅值为$|h|\omega_n/4$；精确主共振$\Omega=2\omega_n$时，存在增长指数$\sigma_0=|h|\omega_n/4+O(h^2\omega_n)$。
(3) 置失谐$\Delta=\Omega-2\omega_n$且$\Delta=O(h\omega_n)$，慢幅本征指数为$\sigma_\pm=\pm\sqrt{(h\omega_n/4)^2-(\Delta/2)^2}$（保留最低非零阶）。因此主不稳定舌满足$|\Delta|<|h|\omega_n/2$；超出该范围，根号转为纯虚，慢变量只作拍振，不出现指数幅度增长。
(4) 弱阻尼使慢幅整体多出$-\beta_n$，故主共振的净增长率$\sigma\simeq|h|\omega_n/4-\beta_n$。当$|h|\omega_n>4\beta_n$才克服耗散，在零失谐时可产生参数不稳定。增长能量来自外部周期改变张力的驱动机构，不违反能量守恒；若振幅不再小，非线性张紧和材料损耗会限制指数增长。
